\chapter{Part 1 --- 1-bit ALU Implementation}
\label{ch:part1}

With the architecture established in \cref{ch:arch}, this chapter takes the
1-bit slice through the full flow: functional simulation, synthesis,
place-and-route, physical verification and timing. The screenshots are the
actual tool captures from the project.

\section{Functional simulation}
Each sub-block is verified in Xcelium with a self-checking testbench that
sweeps \emph{all} input combinations, and the assembled slice is then
verified against the functional table. Because every block is small, the
test vectors are exhaustive rather than random.

\subsection{4:1 multiplexer}
The multiplexer is checked by driving two complementary operand sets and
sweeping the two select bits; the output must equal the selected input
(\cref{tab:mux,fig:muxsim}).

\begin{table}[htbp]
\centering
\caption{4:1 multiplexer truth table (verification vectors).}
\label{tab:mux}
\small
\begin{tabular}{@{}cccc cc c@{}}
\toprule
$a$ & $b$ & $c$ & $d$ & S1 & S0 & $y$\\
\midrule
0&1&0&1 & 0&0 & 0\\
0&1&0&1 & 0&1 & 1\\
0&1&0&1 & 1&0 & 0\\
0&1&0&1 & 1&1 & 1\\
1&0&1&0 & 0&0 & 1\\
1&0&1&0 & 0&1 & 0\\
1&0&1&0 & 1&0 & 1\\
1&0&1&0 & 1&1 & 0\\
\bottomrule
\end{tabular}
\end{table}

\reportfig{0.95}{fig01_mux_sim}{4:1 multiplexer testbench simulation
(Xcelium/SimVision). The output \code{y} tracks the input selected by
\code{S1\,S0}.}{fig:muxsim}

\subsection{Arithmetic circuit}
The arithmetic unit is swept over $a$, $b$, $C_{in}$ and \code{S1\,S0};
the result $d$ and carry $C_{out}$ match the expected
add/subtract/increment/decrement behaviour of \cref{tab:arith}.

\begin{table}[htbp]
\centering
\caption{Arithmetic-circuit truth table. $d$ and $C_{out}$ are the full-adder
outputs for the selected operand ($0$, $B$, $\overline{B}$, $1$).}
\label{tab:arith}
\small
\begin{tabular}{@{}ccc cc cc@{}}
\toprule
$a$ & $b$ & $C_{in}$ & S1 & S0 & $d$ & $C_{out}$\\
\midrule
0&0&0 & 0&0 & 0&0\\
0&0&1 & 0&0 & 1&0\\
0&0&0 & 0&1 & 0&0\\
0&0&1 & 0&1 & 1&0\\
0&0&0 & 1&0 & 1&0\\
0&0&1 & 1&1 & 0&1\\
1&0&0 & 0&0 & 1&0\\
1&0&0 & 0&1 & 1&0\\
1&0&1 & 1&0 & 1&1\\
1&0&0 & 1&1 & 0&1\\
1&1&0 & 0&0 & 1&0\\
1&1&0 & 0&1 & 0&1\\
1&1&1 & 1&0 & 0&1\\
1&1&1 & 1&1 & 1&1\\
\bottomrule
\end{tabular}
\end{table}

\reportfig{0.95}{fig02_arith_sim}{Arithmetic-circuit testbench simulation.
Sum \code{d} and carry-out track the selected arithmetic operation.}
{fig:arithsim}

\subsection{Logic circuit}
Sweeping $a$, $b$ and \code{S1\,S0} confirms the four bitwise functions
(\cref{tab:logic,fig:logicsim}).

\begin{table}[htbp]
\centering
\caption{Logic-circuit truth table: $e=A\!\cdot\!B$, $A\!+\!B$,
$A\!\oplus\!B$, $\overline{A}$ for \code{S1\,S0}$=00,01,10,11$.}
\label{tab:logic}
\small
\begin{tabular}{@{}cc cc c@{}}
\toprule
$a$ & $b$ & S1 & S0 & $e$\\
\midrule
0&0 & 0&0 & 0\\
0&0 & 0&1 & 0\\
0&0 & 1&0 & 0\\
0&0 & 1&1 & 1\\
1&0 & 0&0 & 0\\
1&0 & 0&1 & 1\\
1&0 & 1&0 & 1\\
1&0 & 1&1 & 0\\
1&1 & 0&0 & 1\\
1&1 & 0&1 & 1\\
1&1 & 1&0 & 0\\
1&1 & 1&1 & 0\\
\bottomrule
\end{tabular}
\end{table}

\reportfig{0.95}{fig03_logic_sim}{Logic-circuit testbench simulation. The
output \code{e} realises AND / OR / XOR / NOT under \code{S1\,S0}.}
{fig:logicsim}

\subsection{Assembled 1-bit ALU}
Finally the full slice is exercised across the operations of
\cref{tab:func}. \Cref{fig:alu1sim} shows the waveform; \code{f} and
\code{Cout} follow the specified function for every \sel/\bus{Cin}
combination.

\reportfig{0.95}{fig04_alu1_sim}{1-bit ALU testbench simulation covering the
arithmetic, logic and shift operations of \cref{tab:func}.}{fig:alu1sim}

\section{Synthesis}
The verified RTL is synthesised in Genus against the \texttt{gsclib045}
library at the \texttt{PVT\_0P9V\_125C} corner. Synthesis elaborates the
module hierarchy, maps it onto standard cells, and optimises for the given
constraints. \Cref{fig:alu1synth} shows the resulting gate-level
schematic.

\reportfig{0.9}{fig05_alu1_synth}{Genus synthesis of the 1-bit ALU:
elaborated, technology-mapped gate-level netlist.}{fig:alu1synth}

\subsection{Area}
The slice maps to \textbf{12 cells} totalling \textbf{\SI{39.67}{\micro\meter\squared}}
(\cref{tab:alu1area}, \cref{fig:alu1area}). Note that two \texttt{CLKBUFX20}
drive buffers account for \SI{41.4}{\percent} of the area: these are large
buffers inserted to satisfy the (deliberately heavy) output load used in
the timing setup, not logic proper.

\begin{table}[htbp]
\centering
\caption{1-bit ALU area breakdown by cell type (Genus).}
\label{tab:alu1area}
\small
\begin{tabular}{@{}l r r r@{}}
\toprule
Type & Instances & Area (\umsq) & Share \\
\midrule
Inverter & 1  & 0.684  & \SI{1.7}{\percent}\\
Buffer   & 2  & 16.416 & \SI{41.4}{\percent}\\
Logic    & 9  & 22.572 & \SI{56.9}{\percent}\\
\midrule
\textbf{Total} & \textbf{12} & \textbf{39.672} & \SI{100}{\percent}\\
\bottomrule
\end{tabular}
\end{table}

\reportfig{0.6}{fig06_alu1_area}{Genus area report for the 1-bit ALU.}
{fig:alu1area}

\subsection{Power}
Estimated total power is \textbf{\SI{0.998}{\micro\watt}}, of which
\SI{98.6}{\percent} is dynamic switching power --- expected for a small,
purely combinational block with no clocked elements (\cref{fig:alu1power}).

\reportfig{0.75}{fig07_alu1_power}{Genus power report for the 1-bit ALU
(leakage / internal / switching).}{fig:alu1power}

\section{Place and route}
The netlist is imported into Innovus, floorplanned, placed and routed.
\Cref{fig:alu1layout} shows the placed-and-routed 1-bit slice with its
power rails and signal routing.

\reportfig{0.6}{fig08_alu1_layout}{Innovus placement/routing of the 1-bit
ALU.}{fig:alu1layout}

\section{Physical verification}
Two signoff checks confirm the layout is manufacturable and matches the
netlist:
\begin{itemize}
  \item \textbf{DRC} --- the geometry obeys all \SI{45}{\nano\meter}
        design rules (\cref{fig:alu1drc});
  \item \textbf{Connectivity} --- every net is fully connected with no
        opens or shorts (\cref{fig:alu1conn}).
\end{itemize}

\begin{figure}[htbp]
\centering
\begin{subfigure}{0.48\linewidth}\centering
  \includegraphics[width=\linewidth]{fig09_alu1_drc}
  \caption{DRC results}\label{fig:alu1drc}
\end{subfigure}\hfill
\begin{subfigure}{0.48\linewidth}\centering
  \includegraphics[width=\linewidth]{fig10_alu1_conn}
  \caption{Connectivity check}\label{fig:alu1conn}
\end{subfigure}
\caption{Physical signoff of the 1-bit ALU: (a) design-rule check and
(b) connectivity verification both pass.}
\label{fig:alu1sign}
\end{figure}

\section{Timing and maximum speed}
Post-synthesis static timing analysis reports a critical data-path delay
of \textbf{\SI{3.881}{\nano\second}} through the slice
(\cref{fig:alu1timing}). The path runs from \code{sel\_i[1]} through the
arithmetic multiplexer and full adder to the output. Taking the reciprocal
of the logic delay gives an estimated maximum operating frequency of
\begin{equation}
  f_{\max} \approx \frac{1}{\SI{3.881}{\nano\second}}
           \approx \SI{257}{\mega\hertz},
\end{equation}
with the caveat that this is a relaxed-corner synthesis estimate and that
the reported arrival time includes a large output driver charging the
test load. All paths meet timing with large positive slack against the
(intentionally slow) \SI{1}{\micro\second} analysis clock.

\reportfig{0.85}{fig11_alu1_timing}{Genus timing report for the 1-bit ALU;
the critical path traverses the arithmetic mux and full adder.}
{fig:alu1timing}
